% ============================================================================
%  Capitolo 4 -- Rappresentare l'informazione digitale           (Lezione 04)
% ============================================================================
\chapter{Rappresentare l'informazione digitale}
\label{ch:information}

\epigraphline{\texttt{printf("Hello, world!");}}
  {Brian Kernighan e Dennis Ritchie, \emph{The C Programming Language}, 1978}

\begin{chapgoals}
Alla fine di questo capitolo saprete
\begin{itemize}
\item spiegare la differenza tra informazione analogica e digitale;
\item leggere una configurazione di bit come numero naturale, in binario e in esadecimale;
\item decodificare un colore di \prg come \hex{07E0} nei suoi campi rosso, verde e blu;
\item dire come caratteri e stringhe sono memorizzati, in ASCII e in UTF-8;
\item descrivere come un vettore vive in memoria e come l'assembly ne raggiunge un elemento;
\item distinguere un algoritmo deterministico da uno probabilistico, e collocare tra essi l'apprendimento automatico.
\end{itemize}
\end{chapgoals}
% Generated by tools/textbook/export_listings.js from tools/decks/mastery.js and tools/concepts/registry.json: do not edit.
\begin{chapstatus}{Stato di apprendimento di questo capitolo}
\masterychip{now}{DA COMPRENDERE OGGI}\enspace analogico e digitale; bit, byte, parola; l'esadecimale; la notazione posizionale; un colore è un numero; caratteri e stringhe; l'ordine little-endian; le strutture dati; gli array in memoria; gli algoritmi.\par\smallskip
\masterychip{use}{DA UTILIZZARE OGGI}\enspace shift come moltiplicazione (Capitolo~\ref{ch:integers}) e, usati senza essere studiati nel corso, il programma completo, gli strumenti di PRG32.\par\smallskip
\masterychip{later}{DA COMPRENDERE PIÙ AVANTI}\enspace niente.
\end{chapstatus}


Nel Capitolo~\ref{ch:firstcontact} il lettore ha scritto \hex{07e0} e ha
ottenuto il verde, ha scritto 255 e ha ottenuto una barra che ricomincia, ha
scritto \code{.asciz} e ha ottenuto un saluto, ha scritto \code{.word} e ha
ottenuto un contatore. Questo capitolo dà a quelle cose il loro nome. La sua
tesi sta in una frase, che il resto del capitolo illustra quattro volte:
\emph{nella macchina tutto è una configurazione di bit, e il significato di una
configurazione è una convenzione}. I due capitoli successivi trattano
formalmente i numeri.

% ----------------------------------------------------------------------------
\section{L'informazione}

\subsection{Il flusso dei dati}

Ogni uso di un calcolatore ha la stessa forma. Dei \emph{dati} entrano dal
mondo: un tasto, un sensore, un file. Un \emph{programma}, l'incarnazione di un
algoritmo, li trasforma. Dei \emph{risultati} escono: un'immagine, un suono, un
numero. E qualcuno riceve \emph{conoscenza}, cioè informazione in base alla
quale è possibile agire. Su \prg le quattro fasi sono, nell'ordine, la funzione
\code{prg32\_input\_read}, la \code{update} del programmatore, la \code{draw}
del programmatore, e il giocatore che reagisce allo schermo.

L'ultima fase ha richiesto, finora, una persona. L'intelligenza artificiale, a
cui torna la fine di questo capitolo, è un'intelligenza di un tipo suo proprio:
un potente riconoscimento di schemi, senza autocoscienza e senza esperienza
soggettiva.

\subsection{Informazione analogica e digitale}

\begin{definition}{Analogico e digitale}
Una grandezza \term{analogica} varia con continuità: una tensione, una
pressione sonora, le lancette di un orologio. Una grandezza \term{digitale} è
una sequenza di simboli presi da un alfabeto finito: le cifre di un display.
\end{definition}

\term{Digitalizzare} una grandezza analogica significa farle due cose
(Figura~\ref{fig:info:sampling}): \term{campionarla} nel tempo, leggendone il
valore a istanti regolari, e \term{quantizzarla} nel valore, sostituendo ogni
lettura con il più vicino tra un numero finito di livelli.

\begin{figure}[htb]
\centering
\begin{tikzpicture}[x=0.75cm,y=0.42cm]
  \draw[->,prgink] (0,0) -- (12.6,0) node[right,lbl]{tempo};
  \draw[->,prgink] (0,0) -- (0,8.6) node[above,lbl]{valore};
  \foreach \l in {1,...,7} \draw[prgrule,thin] (0,\l) -- (12.2,\l) node[right,font=\ttfamily\scriptsize,text=prgcomment]{\l};
  \draw[prgaccent,very thick,domain=0:12,smooth,samples=80] plot (\x,{4+2.6*sin(\x*38)+0.9*sin(\x*95)});
  \foreach \x/\y in {0/4,1/6,2/6,3/6,4/6,5/4,6/2,7/1,8/2,9/4,10/4,11/6,12/7}{
    \draw[prgamber,thick] (\x,0) -- (\x,\y); \fill[prgamber] (\x,\y) circle (2.2pt);
  }
  \node[lbl,text=prgaccent] at (3.8,8.4) {segnale analogico};
  \node[lbl,text=prgamber] at (8.2,-0.9) {campioni, quantizzati a 3 bit: 4, 6, 6, 6, 6, 4, 2, 1, \ldots};
\end{tikzpicture}
\caption{Digitalizzare: campionare nel tempo e quantizzare nel valore.}
\label{fig:info:sampling}
\end{figure}

In entrambi i passaggi si perde qualcosa, e si guadagna qualcosa di decisivo.
Un numero finito di livelli può essere \emph{contato}, \emph{copiato
esattamente} e \emph{controllato per rilevare errori}; una grandezza continua
si può soltanto misurare, e ogni copia la degrada. Con due livelli il simbolo è
il \term{bit}, che vale 0 oppure 1. Due livelli sono facili da realizzare in
modo affidabile con i transistor, che sono spenti o accesi
(Capitolo~\ref{ch:electronics}), ed è per questo che i calcolatori sono binari.
L'algebra di Boole (Capitolo~\ref{ch:boole1}) è la matematica dei due livelli.

\subsection{Bit, byte, parola}

I bit si trattano a gruppi di dimensioni standard
(Tabella~\ref{tab:info:units}). Il fatto fondamentale che li riguarda è un
conteggio.

\begin{theorem}{Numero di configurazioni}
Con $n$ bit si possono scrivere esattamente $2^{n}$ configurazioni distinte.
\end{theorem}

\noindent\emph{Dimostrazione.} Per induzione su $n$. Un bit ha $2=2^{1}$
configurazioni. Se $n$ bit hanno $2^{n}$ configurazioni, ciascuna può essere
seguita da uno 0 o da un 1, il che dà $2\cdot 2^{n}=2^{n+1}$ configurazioni di
$n+1$ bit. \hfill$\square$

\begin{table}[htb]
\centering\small
\caption{Unità di informazione.}\label{tab:info:units}
\begin{tabularx}{\textwidth}{@{}llrL@{}}
\toprule
\textbf{Unità} & \textbf{Dimensione} & \textbf{Configurazioni} & \textbf{Dove la si è incontrata}\\
\midrule
bit & 1 cifra binaria & 2 & un pulsante del joystick\\
nibble & 4 bit & 16 & una cifra esadecimale\\
byte & 8 bit & 256 & un carattere ASCII; la larghezza della barra (\code{andi 255})\\
mezza parola & 16 bit & 65\,536 & un colore RGB565\\
parola & 32 bit & 4\,294\,967\,296 & un registro \riscv; un'istruzione; \code{.word}\\
KiB, MiB, GiB & $2^{10}$, $2^{20}$, $2^{30}$ byte & -- & dimensioni di memoria: 1 KiB = 1024 byte\\
\bottomrule
\end{tabularx}
\end{table}

\begin{prgconcept}
Il significato di una configurazione è una convenzione, non una proprietà dei
bit. Che cosa significa una configurazione non c'è scritto dentro: è un accordo
tra chi la scrive e chi la legge.
\end{prgconcept}

La dimensione di una parola dipende dall'architettura: su RV32 una parola è di
32 bit per definizione.

\subsection{Contare in base 2 e in base 16}

La Tabella~\ref{tab:info:hex} elenca le sedici configurazioni di quattro bit
con i loro valori in decimale e le loro cifre esadecimali. Le quattro colonne
di un numero binario pesano 8, 4, 2 e~1. Questa tabella va imparata a memoria:
è la tavola pitagorica di questo libro.

\begin{table}[htb]
\centering\small
\caption{Le sedici configurazioni di quattro bit.}\label{tab:info:hex}
\begin{tabular}{@{}ccc@{\qquad\qquad}ccc@{}}
\toprule
\textbf{Dec} & \textbf{Binario} & \textbf{Esa} & \textbf{Dec} & \textbf{Binario} & \textbf{Esa}\\
\midrule
0 & \texttt{0000} & \texttt{0} & 8 & \texttt{1000} & \texttt{8}\\
1 & \texttt{0001} & \texttt{1} & 9 & \texttt{1001} & \texttt{9}\\
2 & \texttt{0010} & \texttt{2} & 10 & \texttt{1010} & \texttt{A}\\
3 & \texttt{0011} & \texttt{3} & 11 & \texttt{1011} & \texttt{B}\\
4 & \texttt{0100} & \texttt{4} & 12 & \texttt{1100} & \texttt{C}\\
5 & \texttt{0101} & \texttt{5} & 13 & \texttt{1101} & \texttt{D}\\
6 & \texttt{0110} & \texttt{6} & 14 & \texttt{1110} & \texttt{E}\\
7 & \texttt{0111} & \texttt{7} & 15 & \texttt{1111} & \texttt{F}\\
\bottomrule
\end{tabular}
\end{table}

\subsection{La notazione posizionale: che cosa significa \texttt{11011001}}

In un sistema \term{posizionale} il valore di una cifra dipende dalla sua
posizione: ogni posizione pesa una potenza della base. La base dieci usa le
potenze di dieci, la base due le potenze di due, la base sedici le potenze di
sedici. La regola è la stessa; cambia solo la base. Il byte \texttt{11011001}
si legge come segue.

\begin{center}\small
\begin{tabular}{r|*{8}{>{\centering\arraybackslash}p{9mm}|}}
\multicolumn{1}{r}{peso} & \multicolumn{1}{c}{$2^7$} & \multicolumn{1}{c}{$2^6$} & \multicolumn{1}{c}{$2^5$} & \multicolumn{1}{c}{$2^4$} & \multicolumn{1}{c}{$2^3$} & \multicolumn{1}{c}{$2^2$} & \multicolumn{1}{c}{$2^1$} & \multicolumn{1}{c}{$2^0$}\\
\multicolumn{1}{r}{} & \multicolumn{1}{c}{\scriptsize 128} & \multicolumn{1}{c}{\scriptsize 64} & \multicolumn{1}{c}{\scriptsize 32} & \multicolumn{1}{c}{\scriptsize 16} & \multicolumn{1}{c}{\scriptsize 8} & \multicolumn{1}{c}{\scriptsize 4} & \multicolumn{1}{c}{\scriptsize 2} & \multicolumn{1}{c}{\scriptsize 1}\\\cline{2-9}
bit & \textbf{1} & \textbf{1} & 0 & \textbf{1} & \textbf{1} & 0 & 0 & \textbf{1}\\\cline{2-9}
\multicolumn{1}{r}{contributo} & \multicolumn{1}{c}{128} & \multicolumn{1}{c}{64} & \multicolumn{1}{c}{0} & \multicolumn{1}{c}{16} & \multicolumn{1}{c}{8} & \multicolumn{1}{c}{0} & \multicolumn{1}{c}{0} & \multicolumn{1}{c}{1}\\
\end{tabular}
\[128 + 64 + 16 + 8 + 1 = 217\]
\end{center}

Il Capitolo~\ref{ch:integers} dà la formula generale e gli algoritmi di
conversione. Qui lo scopo è soltanto leggere un byte.

\subsection{Perché i programmatori scrivono in esadecimale}

Una cifra esadecimale corrisponde esattamente a quattro bit, perché
$16 = 2^{4}$. Convertire tra binario ed esadecimale è quindi una consultazione
di tabella, cifra per cifra, senza aritmetica:
\[
\underbrace{\texttt{0000}}_{\texttt{0}}\;\underbrace{\texttt{0111}}_{\texttt{7}}\;
\underbrace{\texttt{1110}}_{\texttt{E}}\;\underbrace{\texttt{0000}}_{\texttt{0}}
\qquad\Longrightarrow\qquad \hex{07E0}.
\]
Un registro di 32 bit sono otto cifre esadecimali: \hex{00130313} è leggibile,
trentadue zeri e uni non lo sono. Il prefisso \code{0x} indica l'esadecimale in
assembly e in C, e il debugger di \prgqt mostra i registri in questo modo.

Il valore decimale di \hex{07E0} è $7\cdot 256+14\cdot 16 = 2016$, e non serve
assolutamente a nulla: ciò che conta in questo numero è la \emph{configurazione
di bit}, come mostra la prossima sezione.

\subsection{Un colore è un numero: RGB565}

Su \prg un colore sono sedici bit divisi in tre \term{campi}: cinque bit di
rosso, sei di verde, cinque di blu (Figura~\ref{fig:info:rgb}). Il formato si
chiama RGB565. La costante \code{.equ GREEN, 0x07e0} è la configurazione in cui
il verde è al massimo, 63 su 63, e rosso e blu sono a zero. L'occhio è più
sensibile al verde, ed è per questo che il verde riceve il bit in più.

\begin{figure}[htb]
\centering
\begin{bytefield}[bitwidth=0.05\textwidth,bitheight=8mm]{16}
\bitheader[endianness=big]{0,4,5,10,11,15}\\
\bitbox{5}[bgcolor=prgred!35]{\texttt{00000}} & \bitbox{6}[bgcolor=prggreen!35]{\texttt{111111}} & \bitbox{5}[bgcolor=blue!25]{\texttt{00000}}\\
\bitbox[]{5}{\footnotesize rosso: 5 bit\\0 su 31} & \bitbox[]{6}{\footnotesize verde: 6 bit\\63 su 63} & \bitbox[]{5}{\footnotesize blu: 5 bit\\0 su 31}
\end{bytefield}

\medskip
\begin{tikzpicture}[x=1.55cm]
  \foreach \i/\n/\v/\c in {0/BLACK/0000/{0,0,0}, 1/RED/F800/{248,0,0}, 2/GREEN/07E0/{0,252,0}, 3/BLUE/001F/{0,0,248}, 4/YELLOW/FFE0/{248,252,0}, 5/CYAN/07FF/{0,252,248}, 6/MAGENTA/F81F/{248,0,248}, 7/WHITE/FFFF/{255,255,255}}{
    \definecolor{sw}{RGB}{\c}
    \draw[prgsteel,fill=sw] (\i,0) rectangle ++(0.85,0.6);
    \node[font=\sffamily\scriptsize\bfseries] at (\i+0.42,-0.25) {\n};
    \node[font=\ttfamily\scriptsize,text=prgcomment] at (\i+0.42,-0.6) {0x\v};
  }
\end{tikzpicture}
\caption{RGB565: \hex{07E0} è il verde. Sotto, le costanti di colore di
\fname{prg32.h}.}
\label{fig:info:rgb}
\end{figure}

Sedici bit, tre campi. Gli stessi sedici bit potrebbero essere un numero, due
caratteri o mezza istruzione: il significato sta nel modo in cui il programma
li usa.

\subsection{I numeri: naturali, interi, reali}

Gli stessi numeri sono rappresentati con tre convenzioni diverse, che i due
capitoli successivi sviluppano.

\begin{itemize}
\item I \textbf{numeri naturali} sono la configurazione posizionale stessa:
  $n$ bit coprono l'intervallo da 0 a $2^{n}-1$.
\item Gli \textbf{interi con segno} si scrivono in \term{complemento a due}: un
  registro di 32 bit copre da circa $-2$ miliardi a circa $+2$ miliardi
  (Capitolo~\ref{ch:integers}).
\item I \textbf{numeri reali} si scrivono in \term{virgola fissa} oppure nella
  \term{virgola mobile} dello standard IEEE~754, in cui il valore è
  $(-1)^{s}\times 1,f\times 2^{\,e-\mathit{bias}}$
  (Capitolo~\ref{ch:reals}). La precisione singola usa 32 bit, dà circa 7 cifre
  decimali e arriva a circa $3,4\times 10^{38}$; la precisione doppia usa 64
  bit, dà circa 16 cifre e arriva a circa $1,8\times 10^{308}$.
\end{itemize}

% ----------------------------------------------------------------------------
\section{Il testo}

Anche i caratteri sono numeri.

\subsection{ASCII: un carattere, un byte}

Il codice ASCII (1963) assegna un numero tra 0 e 127, sette bit, a ciascuno di
128 caratteri: le lettere dell'inglese, le cifre, la punteggiatura e alcuni
codici di controllo (Tabella~\ref{tab:info:ascii}).

\begin{table}[htb]
\centering\small
\caption{Alcuni caratteri ASCII.}\label{tab:info:ascii}
\begin{tabularx}{\textwidth}{@{}lrllL@{}}
\toprule
\textbf{Carattere} & \textbf{Decimale} & \textbf{Esa} & \textbf{Binario} & \textbf{Nota}\\
\midrule
\texttt{A} & 65 & \hex{41} & \texttt{0100 0001} & le maiuscole cominciano da 65\\
\texttt{a} & 97 & \hex{61} & \texttt{0110 0001} & come \texttt{A} con un bit in più: $97-65=32$\\
\texttt{0} & 48 & \hex{30} & \texttt{0011 0000} & il carattere zero non è il numero zero\\
spazio & 32 & \hex{20} & \texttt{0010 0000} & un carattere come un altro\\
a capo & 10 & \hex{0A} & \texttt{0000 1010} & un carattere di controllo\\
NUL & 0 & \hex{00} & \texttt{0000 0000} & segna la fine di una stringa C\\
\bottomrule
\end{tabularx}
\end{table}

La tabella è stata progettata con cura. Maiuscole e minuscole differiscono solo
nel bit~5, così che \code{xori} con \hex{20} scambia maiuscola e minuscola: un
primo assaggio di algebra di Boole applicata al testo. E i caratteri delle
cifre sono consecutivi, così che sottraendo 48 dal carattere di una cifra si
ottiene il valore della cifra.

\subsection{Stringhe: una sequenza di caratteri in memoria}

Una \term{stringa} è una sequenza di caratteri memorizzati in byte consecutivi.
L'unica questione è come fa chi legge la stringa a sapere dove finisce, e ci
sono due risposte classiche (Figura~\ref{fig:info:strings}). La \emph{stringa
C} segna la fine con un byte di valore zero. La \emph{stringa Pascal} mette
prima la lunghezza.

\begin{figure}[htb]
\centering\small\ttfamily
\begin{tabular}{@{}r@{\quad}*{5}{|>{\centering\arraybackslash}p{7mm}}|@{\quad}l@{}}
\cline{2-6}
\textsf{Stringa C} & C & I & A & O & \cellcolor{prgred!25}\textbackslash0 & \textsf{\footnotesize la fine è segnata da un byte a zero}\\\cline{2-6}
\multicolumn{1}{r}{} & \multicolumn{1}{c}{\scriptsize 43} & \multicolumn{1}{c}{\scriptsize 49} & \multicolumn{1}{c}{\scriptsize 41} & \multicolumn{1}{c}{\scriptsize 4F} & \multicolumn{1}{c}{\scriptsize 00} & \\[6pt]
\cline{2-6}
\textsf{Stringa Pascal} & \cellcolor{prgamber!25}4 & C & I & A & O & \textsf{\footnotesize prima viene la lunghezza}\\\cline{2-6}
\multicolumn{1}{r}{} & \multicolumn{1}{c}{\scriptsize 04} & \multicolumn{1}{c}{\scriptsize 43} & \multicolumn{1}{c}{\scriptsize 49} & \multicolumn{1}{c}{\scriptsize 41} & \multicolumn{1}{c}{\scriptsize 4F} & \\
\end{tabular}

\bigskip
\begin{tabular}{@{}*{12}{|>{\centering\arraybackslash}p{5.2mm}}|@{}}
\hline
H & E & L & L & O & \textvisiblespace & P & R & G & 3 & 2 & \cellcolor{prgred!25}\textbackslash0\\\hline
\multicolumn{1}{c}{\scriptsize 48} & \multicolumn{1}{c}{\scriptsize 45} & \multicolumn{1}{c}{\scriptsize 4C} & \multicolumn{1}{c}{\scriptsize 4C} & \multicolumn{1}{c}{\scriptsize 4F} & \multicolumn{1}{c}{\scriptsize 20} & \multicolumn{1}{c}{\scriptsize 50} & \multicolumn{1}{c}{\scriptsize 52} & \multicolumn{1}{c}{\scriptsize 47} & \multicolumn{1}{c}{\scriptsize 33} & \multicolumn{1}{c}{\scriptsize 32} & \multicolumn{1}{c}{\scriptsize 00}\\
\end{tabular}
\caption{\normalfont Stringhe in memoria; sotto ogni cella, il codice ASCII in esadecimale.
Sotto: \texttt{.asciz "HELLO PRG32"}.}
\label{fig:info:strings}
\end{figure}

La \code{z} di \code{.asciz} significa \emph{zero-terminated}, terminata da
zero: per gli undici caratteri della stringa del
Capitolo~\ref{ch:firstcontact} l'assemblatore memorizza dodici byte. La
funzione \code{prg32\_gfx\_text8} riceve in \reg{a2} soltanto l'indirizzo del
primo byte, e legge finché non trova lo zero. Se lo zero mancasse, la funzione
continuerebbe a leggere ciò che segue in memoria, qualunque cosa sia.

\subsection{Unicode e UTF-8: ogni sistema di scrittura}

Centoventotto caratteri bastano per l'inglese e per nient'altro.
\term{Unicode} assegna un numero, detto \term{code point}, a ogni carattere di
ogni scrittura; \term{UTF-8} memorizza quel numero in un numero di byte da uno
a quattro (Tabella~\ref{tab:info:utf8}). La codifica è compatibile con ASCII:
un testo ASCII puro è già UTF-8 valido. È la codifica del web e dei file
sorgente che il lettore sta scrivendo.

\begin{table}[htb]
\centering\small
\caption{Code point Unicode e loro codifica UTF-8.}\label{tab:info:utf8}
\begin{tabular}{@{}llll@{}}
\toprule
\textbf{Carattere} & \textbf{Code point} & \textbf{Byte UTF-8} & \textbf{Lunghezza}\\
\midrule
A & U+0041 & \texttt{41} & 1 byte: identico ad ASCII\\
\`e & U+00E8 & \texttt{C3 A8} & 2 byte\\
\texteuro & U+20AC & \texttt{E2 82 AC} & 3 byte\\
l'ideogramma per «sole», «giorno» & U+65E5 & \texttt{E6 97 A5} & 3 byte\\
l'emoji della «faccina sorridente» & U+1F600 & \texttt{F0 9F 98 80} & 4 byte\\
\bottomrule
\end{tabular}
\end{table}

\begin{prgwarn}
In UTF-8 la lunghezza di una stringa in byte non è la sua lunghezza in
caratteri.
\end{prgwarn}

\subsection{L'ordine dei byte: little-endian}

Una parola sono quattro byte, e la memoria è indirizzata byte per byte. In che
ordine vanno i quattro byte di una parola? La direttiva
\code{.word 0x12345678} si può memorizzare in due modi
(Figura~\ref{fig:info:endian}).

\begin{figure}[htb]
\centering\small
\begin{tabular}{@{}r@{\quad}*{4}{|>{\centering\arraybackslash\ttfamily}p{13mm}}|@{}}
\multicolumn{1}{r}{} & \multicolumn{1}{c}{\scriptsize indirizzo $A$} & \multicolumn{1}{c}{\scriptsize $A+1$} & \multicolumn{1}{c}{\scriptsize $A+2$} & \multicolumn{1}{c}{\scriptsize $A+3$}\\\cline{2-5}
\textbf{little-endian} (\riscv, x86, ARM) & \cellcolor{prggreen!25}78 & \cellcolor{prgviolet!20}56 & \cellcolor{prgamber!25}34 & \cellcolor{prgaccent!20}12\\\cline{2-5}
\multicolumn{5}{c}{}\\[-6pt]\cline{2-5}
\textbf{big-endian} (protocolli di rete) & \cellcolor{prgaccent!20}12 & \cellcolor{prgamber!25}34 & \cellcolor{prgviolet!20}56 & \cellcolor{prggreen!25}78\\\cline{2-5}
\end{tabular}
\caption{I due ordini dei byte della parola \hex{12345678}.}
\label{fig:info:endian}
\end{figure}

\begin{definition}{Little-endian e big-endian}
Nell'ordine \term{little-endian} il byte meno significativo di una parola è
memorizzato all'indirizzo più basso; nell'ordine \term{big-endian}, il byte più
significativo.
\end{definition}

\riscv è little-endian: è ciò che fanno \ins{lw} e \ins{sw}. L'ordine dei byte
si nota in due circostanze: quando si esamina la memoria in un debugger, e
quando si scambiano dati con un'altra macchina. I nomi vengono dai \emph{Viaggi
di Gulliver}, dove due popoli si fanno la guerra per decidere da quale estremità
vada rotto un uovo.

% ----------------------------------------------------------------------------
\section{Strutture dati}

Dai singoli valori si passa alle collezioni organizzate. Cinque modi di
organizzare i dati ricorrono in tutta l'informatica.

\begin{description}[style=nextline,leftmargin=1.5em]
\item[Vettore (array).] Elementi dello stesso tipo, uno dopo l'altro, raggiunti
  tramite un indice: 32, 11, 34, 22.
\item[Lista.] Elementi collegati ciascuno al successivo: economica per inserire
  e togliere, lenta da raggiungere per posizione.
\item[Pila (stack).] L'ultimo a entrare è il primo a uscire: \emph{push}
  aggiunge in cima, \emph{pop} toglie dalla cima. Il frame che ha salvato
  \reg{ra} nel Capitolo~\ref{ch:firstcontact} stava su uno stack.
\item[Dizionario.] Coppie di chiave e valore: \code{NA} dà \code{NAPOLI},
  \code{PA} dà \code{PALERMO}, \code{MI} dà \code{MILANO}.
\item[Albero.] Una gerarchia: un animale cammina o nuota, ha polmoni o
  branchie, due zampe o quattro.
\end{description}

Nella macchina sono tutti la stessa cosa: byte in memoria, più una regola per
trovare l'elemento successivo.

\subsection{Un vettore in memoria, un elemento in un registro}

Per il vettore la regola è aritmetica (Listato~\ref{lst:info:array}). Un
vettore è un indirizzo di base seguito da elementi di uguale dimensione;
l'elemento $i$ di un vettore di parole sta all'indirizzo $\mathit{base}+4i$,
perché una parola sono quattro byte.

\asmfile[caption={Leggere \code{scores[2]}.},label=lst:info:array]{code/l04/c03.S}

L'istruzione \ins{slli} fa scorrere a sinistra i bit di un registro; uno
scorrimento di due posizioni moltiplica per quattro, come in decimale uno
scorrimento di una posizione moltiplica per dieci. (La proprietà è dimostrata
nel Capitolo~\ref{ch:integers}.) L'espressione \code{scores[i]} di un
linguaggio di alto livello è esattamente queste quattro istruzioni.

\begin{prgwarn}
Che cosa legge \ins{lw} se $i=4$? Ciò che segue il vettore in memoria,
qualunque cosa sia. La macchina non controlla i limiti: lo fa il programmatore.
\end{prgwarn}

% ----------------------------------------------------------------------------
\section{Dati e algoritmi}

Il capitolo si chiude allargando lo sguardo, dai bit di un programma ai dati
del mondo e alle procedure che lavorano su di essi.

\subsection{Dati, big data e data science}

\begin{definition}{Big data}
I \term{big data} sono dati che superano i limiti degli strumenti tradizionali
delle basi di dati.
\end{definition}

Sono strutturati e non strutturati: l'economia, la geolocalizzazione, le serie
temporali, le reti sociali. La loro dimensione si misura in petabyte e
zettabyte, cioè in miliardi di terabyte. Vengono da quattro tipi di sorgente: i
\emph{dati da sensori} (l'ambiente, le posizioni GPS, l'Internet delle cose:
ogni dispositivo misura qualcosa); le \emph{transazioni} (gli acquisti al
supermercato, le carte di credito, gli ordini online); i \emph{dati pubblici}
(atti amministrativi, territorio, consumi: gli open data delle istituzioni); e
i \emph{social media} (preferenze, politica, opinioni: ciò che le persone
scrivono e condividono). La \term{data science} è l'analisi di grandi quantità
di informazione, e le tecnologie che estraggono conoscenza e valore dai dati
sono il calcolo parallelo e le GPU.

\begin{definition}{Data mining}
Il \term{data mining} è l'insieme delle tecniche e delle metodologie per
estrarre informazione o conoscenza da grandi quantità di dati.
\end{definition}

Tra le sue tecniche ci sono gli alberi di decisione e la classificazione,
l'apprendimento di regole e i metodi bayesiani, le reti neurali e il
riconoscimento di schemi.

\subsection{GPGPU: il calcolo generico sui processori grafici}

Un processore grafico (GPU) ha migliaia di core semplici costruiti per fare la
stessa operazione su molti dati insieme. La somma di $N$ elementi, che richiede
$N$ passi uno dopo l'altro, ha bisogno di soli $\log_2 N$ passi paralleli
(Figura~\ref{fig:info:gpgpu}): gli elementi si sommano a coppie, poi le somme
parziali a coppie, e così via. Lo stesso hardware che disegna i videogiochi
oggi addestra le reti neurali. Sotto ogni core ci sono sempre registri,
istruzioni e bit; il Capitolo~\ref{ch:beyond} torna sull'argomento.

\begin{figure}[htb]
\centering
\includegraphics[width=0.6\textwidth]{gpgpu_sum.png}
\caption{Una somma parallela sui core di una GPU: $\log_2 N$ passi.}
\label{fig:info:gpgpu}
\end{figure}

\subsection{Algoritmi e calcolabilità}

Un problema è \term{calcolabile} se può essere risolto da un algoritmo. La
\emph{macchina di Turing} è un calcolatore ideale con un programma, espresso
matematicamente tramite insiemi, relazioni e funzioni. La \emph{macchina di von
Neumann} è l'architettura di calcolo capace di eseguire un programma che
realizza un algoritmo per la macchina di Turing. Un problema che nessun
algoritmo su una macchina di Turing può risolvere non può essere risolto da
alcun calcolatore: questo enunciato, nella sua versione informale, è la tesi di
Church--Turing. Non ogni problema ha un algoritmo; l'esempio classico di
problema non calcolabile è il \emph{problema della fermata}, quello di decidere
se un programma arbitrario prima o poi si fermerà.

\begin{definition}{Algoritmo}
Un algoritmo è una procedura con cinque proprietà. È \textbf{atomico}: i suoi
passi sono elementari e non ulteriormente scomponibili. È \textbf{non
ambiguo}: ogni passo è interpretato in un solo modo diretto dall'esecutore,
umano o macchina. È \textbf{finito}: ha un numero finito di passi e una
quantità finita di dati in ingresso. È \textbf{terminante}: la sua esecuzione
finisce dopo un tempo finito. È \textbf{effettivo}: la sua esecuzione porta a
un risultato unico e ben definito.
\end{definition}

Un'istruzione \riscv è il modello di un passo atomico e non ambiguo: il suo
significato è fissato dalla specifica, bit per bit.

\subsection{Algoritmi deterministici e non deterministici}

Un algoritmo è \term{deterministico} se per ogni istruzione esiste, a parità di
ingresso, esattamente un passo successivo: il percorso dell'algoritmo esiste ed
è unico, $y=f(x)$. Convertire un'unità di misura è deterministico; lo è anche
\code{first\_update}, che somma sempre uno.

Un algoritmo è \term{non deterministico} se, a parità di ingresso, può
comportarsi in modo diverso a ogni esecuzione. Gli algoritmi
\emph{probabilistici} usano un generatore di numeri casuali; gli algoritmi
\emph{concorrenti} possono dipendere dai tempi (le cosiddette race condition).
Un gioco che colloca un nemico chiamando \code{prg32\_random\_number} è non
deterministico.

Il Listato~\ref{lst:info:prob} mostra due algoritmi probabilistici in
pseudo-codice. Il primo verifica se un numero è primo: la risposta «composto» è
sempre giusta, la risposta «probabilmente primo» è giusta con probabilità
schiacciante. Scambiare una minuscola possibilità di errore con un enorme
guadagno di velocità è il modo in cui la crittografia reale verifica i numeri
primi. (Il test di Fermat mostrato è il più semplice; le librerie reali usano
il test di Miller--Rabin, che ne chiude le falle.) Il secondo è una passeggiata
aleatoria, che dà un risultato diverso a ogni esecuzione: le simulazioni ne
fanno la media su molte.

\txtfile[caption={Due algoritmi probabilistici.},label=lst:info:prob]{code/l04/c04.txt}

\subsection{L'apprendimento automatico}

L'\term{apprendimento automatico} (machine learning) è un metodo di analisi dei
dati che automatizza la costruzione di un modello analitico: il calcolatore
trova schemi nascosti senza essere stato esplicitamente programmato su dove
cercare. In un programma classico la funzione $y=f(x)$ è scritta da una
persona; nell'apprendimento automatico $f$ viene appresa dagli esempi. I
modelli generativi di oggi sono la stessa idea su scala enorme: un grande
modello linguistico predice il prossimo pezzo di testo a partire da miliardi di
numeri appresi.

Quei numeri sono bit in memoria, moltiplicati e sommati da processori:
l'argomento di questo libro. E un modello può sbagliare con grande sicurezza,
quindi per controllarlo servono i fondamenti. È la ragione della regola
enunciata nel Capitolo~\ref{ch:intro}: l'esame verifica che i fondamenti stiano
nella testa dello studente.

% ----------------------------------------------------------------------------
% Generated by tools/textbook/export_listings.js from decks/plan.js and code/l04.S: do not edit.
\section{Il programma di questo capitolo}\label{sec:lab:l04}

I listati di questo capitolo sono parti di un solo programma,
\fname{code/l04.S} nel repository del corso, stampato per intero nel
Listato~\ref{lst:lab:l04}. I test automatici del materiale del corso lo
compilano e lo eseguono sull'esecutore senza interfaccia di \prgqt{}, con
ingressi programmati, e controllano il frame; la
Figura~\ref{fig:lab:l04} è il frame scritto dal test.

\begin{description}[style=nextline,leftmargin=1.5em]
\item[Sullo schermo.] Quattro barre orizzontali, una per elemento dell'array (32, 11, 34, 22, quattro pixel per unità); l'elemento 2, quello che il listato carica, è giallo.
\item[Che cosa premere.] Niente: il programma va da solo.
\end{description}

\begin{figure}[htb]
\centering
\includegraphics[width=0.5\textwidth]{lab_l04.png}
\caption{\fname{l04.S} su \prgqt{}.}\label{fig:lab:l04}
\end{figure}

\asmfile[caption={Il programma completo del capitolo, \fname{code/l04.S}.},label=lst:lab:l04]{code/l04/program.S}

Per eseguirlo su \prgqt{} i comandi sono quelli del
Listato~\ref{lst:lab:l04:run}; il Paragrafo~\ref{sec:prg:testing} dà gli
stessi passi per QEMU e per la scheda.

\begin{shell}[caption={Compilare \fname{l04.S} ed eseguirlo su \prgqt{}.},label=lst:lab:l04:run]
# 1. start PRG32-QT on this computer (it listens on port 8080)

# 2. in the PRG32 folder, with ESP-IDF active: build the cartridge
python3 -m prg32 cartridge build ~/arc12/code/l04.S --portable \
    --entry-prefix l04 --name l04 --out l04.prg32

# 3. send it to PRG32-QT and start it
python3 -m prg32 esp32c6 upload-and-run l04.prg32 \
    --url http://localhost:8080 --slot cart0

# 4. look, press, compare with the previous slide
\end{shell}


% ----------------------------------------------------------------------------
% Generated by tools/textbook/export_listings.js from tools/concepts/assessments.json: do not edit.
\begin{chapstatus}{Il filo del capitolo}
I bit da soli non significano nulla: un codice dà loro un significato, e gli stessi 32 bit possono essere un numero, un colore, quattro lettere o un'istruzione. Prossimo passo: i codici degli interi, e che cosa succede quando i bit finiscono.
\par\medskip\textbf{Biglietto d'uscita} (da due a quattro minuti, da soli, su carta):
\begin{enumerate}[nosep,leftmargin=2em]
\item Un termometro a mercurio e un display che mostra 21,5 misurano entrambi la temperatura. Quale è analogico e quale digitale? Perché la lettura digitale si può copiare mille volte senza peggiorare?
\item Il pulsante Annulla di un editor cancella prima la modifica più recente, poi quella precedente. Quale delle cinque strutture dati è adatta, e perché non un array letto per indice?
\end{enumerate}
\end{chapstatus}


\section{Verificate di aver capito}

\begin{quiz}{Quanti valori diversi può rappresentare un byte?}
\item 8
\item 16
\item 255
\item 256
\end{quiz}

\begin{quiz}{Il colore \prg \hex{F81F} è:}
\item rosso
\item magenta
\item ciano
\item giallo
\end{quiz}

\begin{quiz}{Quanti byte occupa in memoria \code{.asciz "CIAO"}?}
\item 3
\item 4
\item 5
\item 8
\end{quiz}

\begin{quiz}{Su \riscv, \code{.word 0x0000ABCD} è memorizzata all'indirizzo $A$. Quale byte si trova all'indirizzo $A$?}
\item \hex{00}
\item \hex{AB}
\item \hex{CD}
\item \hex{DC}
\end{quiz}

\section{Esercizi}

\begin{exercise}[su \prg]
\begin{enumerate}[nosep]
\item Cambiate la barra del Capitolo~\ref{ch:firstcontact} in magenta, poi in
  un grigio scuro a vostra scelta: calcolate prima a mano il valore RGB565.
\item Scrivete il vostro nome con \code{prg32\_gfx\_text8}. Quanti byte occupa
  la stringa?
\item Dichiarate un vettore di quattro colori con \code{.word} e disegnate
  quattro barre, una per colore, leggendo ogni colore con \ins{lw}.
\item In \prgqt aprite la vista della memoria sulla vostra stringa e sul vostro
  vettore: trovate lo zero finale e verificate l'ordine dei byte.
\end{enumerate}
\expected{Quattro barre in quattro colori; nella vista della memoria la stringa
finisce con \texttt{00} e ogni \code{.word} mostra per primo il byte meno
significativo. Un grigio scuro, per confronto: rosso 8 su 31, verde 16 su 63,
blu 8 su 31 è \texttt{01000\,010000\,01000}, cioè \hex{4208}.}
\end{exercise}

\section{Glossario}

\begin{glossarytable}
bit / byte / word & bit / byte / parola\\
positional notation & notazione posizionale\\
hexadecimal & esadecimale\\
sampling & campionamento\\
quantisation & quantizzazione\\
byte order (endianness) & ordinamento dei byte\\
string & stringa\\
array & vettore (array)\\
linked list & lista concatenata\\
stack & pila\\
tree & albero\\
machine learning & apprendimento automatico\\
\end{glossarytable}

\section{Conclusioni e letture}

L'informazione digitale è una configurazione di bit più una convenzione per
leggerla. Gli stessi sedici bit possono essere un numero, un colore o due
caratteri: lo decide il programma. L'esadecimale è una scrittura compatta del
binario, quattro bit per cifra. Il testo è fatto di numeri (ASCII, UTF-8); le
strutture sono byte più una regola per trovare l'elemento successivo. I dati
sono la moneta di questo regno, e gli algoritmi trasformano i dati in
informazione adatta a decidere.

\paragraph{Letture.} \citet{patterson2020}, paragrafi 2.4 e~2.9; Charles
Petzold, \emph{Code}, i capitoli su bit, byte e testo; per l'apprendimento
automatico, Sebastian Raschka, \emph{Machine Learning with Python}.